% ----------
% Plantilla LaTeX para Reportes de Proyectos
% Adaptada para: Sesión 1 — Proyecto Final Bases de Datos
% Universidad Nacional de Colombia - Sede Bogotá
% ----------
\documentclass[12pt, letterpaper, twoside]{article}

% --- PAQUETES PRINCIPALES ---
\usepackage{geometry}
\geometry{top=2.5cm, bottom=2.5cm, left=3cm, right=2.5cm, headheight=15pt}
\usepackage[utf8]{inputenc}
\usepackage[T1]{fontenc}
\usepackage[spanish, es-tabla, es-nodecimaldot]{babel}
\usepackage{booktabs}
\usepackage{longtable}
\usepackage{fancyhdr}
\usepackage{helvet}
\renewcommand{\familydefault}{\sfdefault}
\usepackage{csquotes}
\usepackage{graphicx}
\usepackage{parskip}
\usepackage{titling}
\usepackage{etoolbox}
\usepackage{xcolor}
\usepackage{titlesec}
\usepackage{enumitem}
\usepackage{microtype}
\usepackage{array}

% ----------
% DEFINICIÓN DE COLORES INSTITUCIONALES
% ----------
\definecolor{azulUNAL}{RGB}{0, 57, 115}

% ----------
% CONFIGURACIÓN DE SECCIONES
% ----------
\titleformat{\section}{\large\bfseries\color{azulUNAL}}{\thesection.}{0.5em}{}[\titlerule]
\titleformat{\subsection}{\normalsize\bfseries}{\thesubsection.}{0.5em}{}

% ----------
% COMANDOS PERSONALIZADOS
% ----------
\newcommand{\runningtitle}[1]{\def\therunningtitle{#1}}
\newcommand{\runningauthor}[1]{\def\therunningauthor{#1}}
\newcommand{\affiliation}[1]{\def\theaffiliation{#1}}
\newcommand{\department}[1]{\def\thedepartment{#1}}
\newcommand{\memoid}[1]{\def\thememoid{#1}}
\newcommand{\theyear}[1]{\def\thetheyear{#1}}
\newcommand{\mydate}[1]{\def\themydate{#1}}

% ----------
% ENCABEZADO Y PIE DE PÁGINA
% ----------
\fancyhf{}
\fancyhead[LE,RO]{\small\thepage}
\fancyhead[LO]{\small\nouppercase{\therunningtitle}}
\fancyhead[RE]{\small\nouppercase{\therunningauthor}}
\fancyfoot[C]{\small Universidad Nacional de Colombia}
\renewcommand{\headrulewidth}{0.4pt}
\renewcommand{\footrulewidth}{0.4pt}
\pagestyle{fancy}

% --- Estilo de la primera página ---
\fancypagestyle{firstpage}{%
  \fancyhf{}
  \fancyfoot[C]{\small Universidad Nacional de Colombia}
  \renewcommand{\headrulewidth}{0pt}
  \renewcommand{\footrulewidth}{0.4pt}
}

% ----------
% RESUMEN EN ESPAÑOL
% ----------
\renewenvironment{abstract}{%
  \small\begin{center}%
  {\bfseries\abstractname\vspace{-.5em}}%
  \end{center}%
  \quotation}{%
  \endquotation}

% ----------
% PORTADA
% ----------
\preauthor{\begin{center}\large}\postauthor{\end{center}}
\predate{\begin{center}\large}\postdate{\end{center}}

\renewcommand{\maketitle}{%
\begin{titlepage}
  \centering
  \vspace*{1.5cm}
  {\large\bfseries\color{azulUNAL}
    Universidad Nacional de Colombia\\[0.3em]
    \normalsize\color{black}\theaffiliation\\[0.1em]
    \thedepartment\par}
  \vspace{1cm}
  \rule{\linewidth}{1.5pt}\\[0.4em]
  {\LARGE\bfseries\thetitle\par}
  \rule{\linewidth}{1.5pt}
  \vspace{1cm}

  % Profesor
  {\normalsize\textbf{Profesor}\\[0.2em]
  Cristhian Fernando Lara Espejo\par}

  \vspace{1.2cm}

  % Integrantes
  {\normalsize\textbf{Nombres y códigos de los estudiantes:}\\[0.4em]
  Santiago Rafael Puentes Durán: 1032936643\\[0.2em]
  Sergio Gabriel Chaves Mosquera: 1014990992\\[0.2em]
  Cristian Hernández Muñoz: 1141320791\\[0.2em]
  Cesar Samuel Sánchez Romero: 1011100422\\[0.2em]
  Edison Stiven Quintero Motta: 1028561323\par}

  \vfill
  {\small\color{gray}\thetheyear}
\end{titlepage}
}

% ----------
% VARIABLES DEL DOCUMENTO
% ----------
\title{\textbf{Proyecto Final --- Bases de Datos\\[0.3em]
Sesión 1: Kickoff y Diseño Conceptual\\[0.3em]
\guillemotleft UniGrades\guillemotright}}
\runningtitle{UniGrades --- Sesión 1}
\runningauthor{Puentes, Chaves, Hernández, Sánchez, Quintero}
\affiliation{Sede Bogotá}
\department{Curso: Bases de Datos}
\memoid{}
\theyear{2026}
\mydate{16 de julio de 2026}

% ==========================================================
%                    CUERPO DEL DOCUMENTO
% ==========================================================
\begin{document}
\maketitle
\thispagestyle{firstpage}
\newpage

% --- RESUMEN ---
\begin{abstract}
\noindent
Este documento presenta el diseño conceptual de \textbf{UniGrades}, una aplicación
web para la gestión del progreso académico universitario. Se define el dominio del
problema, se identifican las reglas de negocio que rigen el sistema, se construye
el diagrama Entidad-Relación con sus 9 entidades, relaciones de cardinalidad y
participación, y se elabora el diccionario de datos preliminar. El modelo soporta
múltiples universidades y programas curriculares, permitiendo al estudiante registrar
materias, configurar componentes de evaluación acumulativos y calcular promedios
ponderados por créditos por semestre y de forma global.
\end{abstract}

\vspace{1.5cm}
\tableofcontents
\newpage

% ==========================================================
\section{Descripción del Dominio}
\label{sec:dominio}
% ==========================================================

UniGrades es una aplicación web orientada a estudiantes universitarios que centraliza
el seguimiento de su progreso académico a lo largo de toda la carrera. El sistema
permite registrar las materias cursadas en cada período académico, configurar los
componentes de evaluación de cada asignatura (parciales, talleres, proyectos, etc.)
y registrar las calificaciones obtenidas, calculando automáticamente promedios por
materia, por semestre y el promedio global acumulado.

A diferencia de los sistemas institucionales (como el SIA de la UNAL), UniGrades
es una herramienta personal del estudiante: él mismo configura su plan de trabajo,
registra sus notas a medida que las recibe y visualiza su avance por tipología de
créditos. El sistema está diseñado para ser agnóstico a la institución, soportando
cualquier universidad, programa y escala de calificación.

Una base de datos relacional es indispensable porque el dominio involucra múltiples
entidades fuertemente relacionadas (universidad, programa, tipología, materia,
usuario, semestre, inscripción, componente y nota), datos históricos transaccionales
(semestres cursados, notas registradas por fecha), restricciones de integridad no
triviales (los porcentajes de un componente deben sumar 100, la nota mínima de
aprobación varía por materia) y consultas analíticas que cruzan varias dimensiones
(promedio ponderado por créditos, avance por tipología, \textit{roadmap} sugerido
frente al real).

% ==========================================================
\section{Reglas de Negocio}
\label{sec:reglas}
% ==========================================================

Las siguientes reglas describen las restricciones y comportamientos obligatorios
del sistema. Cada regla indica las entidades involucradas y la restricción que
impone al modelo relacional.

\begin{enumerate}[leftmargin=*, label=\textbf{RN-\arabic*.}]

  \item \textbf{Un estudiante pertenece a exactamente un programa.}\\
  \textit{Entidades:} \texttt{usuario}, \texttt{programa}.\\
  Un \texttt{usuario} debe tener asociado un único \texttt{programa} curricular al
  momento del registro. No puede existir un usuario sin programa asignado
  (\texttt{usuario\_programa\_id NOT NULL}).

  \item \textbf{Los porcentajes de los componentes de una materia deben sumar 100.}\\
  \textit{Entidades:} \texttt{componente}, \texttt{materia\_usuario}.\\
  La suma de \texttt{componente\_porcentaje} de todos los componentes asociados a
  una misma \texttt{materia\_usuario} debe ser exactamente 100. Esta restricción
  requiere validación a nivel de aplicación o \textit{trigger}, ya que no puede
  expresarse con \texttt{NOT NULL} ni \texttt{UNIQUE}.

  \item \textbf{Un estudiante no puede inscribir la misma materia dos veces en el mismo semestre.}\\
  \textit{Entidades:} \texttt{materia\_usuario}, \texttt{semestre}, \texttt{materia}.\\
  La combinación (\texttt{materia\_usuario\_usuario\_id},
  \texttt{materia\_usuario\_materia\_id}, \texttt{materia\_usuario\_semestre\_id})
  debe ser única. Un mismo estudiante puede cursar la misma materia en semestres
  distintos (habilitación o repetición), pero no dos veces en el mismo período.

  \item \textbf{El estado de una materia sigue transiciones definidas.}\\
  \textit{Entidades:} \texttt{materia\_usuario}.\\
  El atributo \texttt{materia\_usuario\_estado} solo puede tomar los valores
  \texttt{en\_curso}, \texttt{aprobada}, \texttt{reprobada} o \texttt{retirada}.
  Una materia con estado \texttt{aprobada} o \texttt{reprobada} no puede volver a
  \texttt{en\_curso} en el mismo registro; debe crearse una nueva inscripción en
  otro semestre.

  \item \textbf{El valor de una nota debe estar dentro del rango válido de la escala.}\\
  \textit{Entidades:} \texttt{nota}.\\
  \texttt{nota\_valor} debe estar entre 0.0 y 5.0 (escala colombiana estándar).
  Valores fuera de rango son rechazados mediante una restricción \texttt{CHECK}.

  \item \textbf{Las tipologías de nivelación no cuentan en el promedio acumulado.}\\
  \textit{Entidades:} \texttt{tipologia}, \texttt{materia}, \texttt{nota}.\\
  Las materias cuya tipología tiene \texttt{tipologia\_cuenta\_promedio = 0}
  (p.\,ej.\ Inglés I--IV en la UNAL) se excluyen del cálculo del promedio global
  y del promedio por semestre, aunque sí se registran y validan en el sistema.

  \item \textbf{Un componente acumula múltiples notas cuyo promedio define su calificación.}\\
  \textit{Entidades:} \texttt{componente}, \texttt{nota}.\\
  Un \texttt{componente} puede tener cero o más registros de \texttt{nota}
  (relación 1:N). La calificación del componente se calcula como el promedio
  aritmético de todas las notas registradas en ese momento. Esto permite modelar
  componentes acumulativos (p.\,ej.\ \textit{Talleres -- 20\,\%}) donde la
  cantidad de entregas no se conoce de antemano.

\end{enumerate}

% ==========================================================
\section{Diagrama Entidad-Relación}
\label{sec:er}
% ==========================================================

El diagrama E-R fue elaborado en draw.io siguiendo la notación crow's foot.
Se incluye como archivo adjunto \texttt{Diagrama E-R.drawio} y su versión
exportada \texttt{Diagrama E-R.drawio.pdf} en el repositorio del proyecto
(\url{https://github.com/sechaves/UniGrades}).

\subsection{Entidades del modelo}

El modelo cuenta con \textbf{9 entidades base}, entre maestras y transaccionales:

\begin{itemize}[leftmargin=1.5em]
  \item \textbf{UNIVERSIDAD} --- maestra; almacena las instituciones soportadas.
  \item \textbf{PROGRAMA} --- maestra; carrera curricular dentro de una universidad.
  \item \textbf{TIPOLOGIA} --- maestra; categoría de materia dentro de un programa
        (Obligatoria, Optativa, Libre Elección, Nivelación, etc.).
  \item \textbf{MATERIA} --- maestra; asignatura del programa con semestre sugerido (\textit{roadmap}).
  \item \textbf{USUARIO} --- maestra; el estudiante registrado en la plataforma.
  \item \textbf{SEMESTRE} --- transaccional; período académico real cursado por el estudiante.
  \item \textbf{MATERIA\_USUARIO} --- asociativa ternaria (M:N con atributos propios);
        vincula usuario, materia y semestre. Registra el estado de la materia.
  \item \textbf{COMPONENTE} --- transaccional; evaluación configurada por el estudiante
        (nombre, porcentaje, nota mínima).
  \item \textbf{NOTA} --- transaccional; calificación individual dentro de un componente.
\end{itemize}

\subsection{Relaciones M:N con atributos propios}

\begin{enumerate}[leftmargin=1.5em]
  \item \textbf{MATERIA\_USUARIO} (USUARIO $\times$ MATERIA $\times$ SEMESTRE): tabla
  asociativa ternaria. Atributo propio: \texttt{materia\_usuario\_estado}
  (\texttt{en\_curso / aprobada / reprobada / retirada}).
  \item \textbf{COMPONENTE} como extensión de MATERIA\_USUARIO: cada inscripción
  genera sus propios componentes de evaluación con atributos propios
  \texttt{componente\_porcentaje} y \texttt{componente\_nota\_minima}.
\end{enumerate}

\subsection{Jerarquía y datos temporales}

La entidad \textbf{TIPOLOGIA} incorpora el atributo \texttt{tipologia\_cuenta\_promedio},
que establece una jerarquía funcional entre categorías: aquellas con valor 0 son de
nivelación y se excluyen de los cálculos de promedio. Adicionalmente, la entidad
\textbf{PROGRAMA} articula de forma jerárquica a \textbf{TIPOLOGIA} y \textbf{MATERIA},
modelando la estructura curricular: universidad $\rightarrow$ programa $\rightarrow$
tipología $\rightarrow$ materia.

La entidad \textbf{SEMESTRE} registra \texttt{semestre\_year} y \texttt{semestre\_periodo},
habilitando consultas históricas por rango de períodos académicos. La entidad
\textbf{NOTA} registra \texttt{nota\_fecha\_registro}, permitiendo el seguimiento
cronológico del progreso del estudiante.

\subsection{Resumen de relaciones y cardinalidades}

\begin{center}
\begin{tabular}{llll}
\toprule
\textbf{Entidad origen} & \textbf{Cardinalidad} & \textbf{Entidad destino} & \textbf{Participación destino} \\
\midrule
UNIVERSIDAD      & 1 : N & PROGRAMA         & Total \\
PROGRAMA         & 1 : N & TIPOLOGIA        & Total \\
PROGRAMA         & 1 : N & MATERIA          & Total \\
TIPOLOGIA        & 1 : N & MATERIA          & Total \\
PROGRAMA         & 1 : N & USUARIO          & Total \\
USUARIO          & 1 : N & SEMESTRE         & Total \\
USUARIO          & 1 : N & MATERIA\_USUARIO  & Total \\
MATERIA          & 1 : N & MATERIA\_USUARIO  & Total \\
SEMESTRE         & 1 : N & MATERIA\_USUARIO  & Total \\
MATERIA\_USUARIO & 1 : N & COMPONENTE       & Total \\
COMPONENTE       & 1 : N & NOTA             & Parcial \\
\bottomrule
\end{tabular}
\end{center}

% ==========================================================
\section{Diccionario de Datos Preliminar}
\label{sec:diccionario}
% ==========================================================

Se documenta cada entidad con sus atributos, tipo de dato sugerido,
restricciones aplicadas y descripción funcional.

\medskip
\noindent\textbf{Nota:} CP = Clave Primaria \quad CE = Clave Externa (Foránea)

% ----------------------------------------------------------
\subsection{Tabla: UNIVERSIDAD}

\begin{center}
\begin{tabular}{p{4cm} p{2.8cm} p{2.8cm} p{4cm}}
\toprule
\textbf{Atributo} & \textbf{Tipo} & \textbf{Restricción} & \textbf{Descripción} \\
\midrule
universidad\_id          & INT            & CP, NO NULO   & Identificador único de la universidad. \\
universidad\_nombre      & VARCHAR(150)   & NO NULO       & Nombre oficial de la institución. \\
universidad\_sigla       & VARCHAR(20)    & Admite nulo   & Sigla o acrónimo (ej.\ UNAL). \\
universidad\_pais        & VARCHAR(80)    & NO NULO       & País sede de la institución. \\
universidad\_ciudad      & VARCHAR(80)    & Admite nulo   & Ciudad sede. \\
universidad\_logo\_url   & VARCHAR(255)   & Admite nulo   & URL del logotipo para la interfaz. \\
universidad\_created\_at & TIMESTAMP      & DEF NOW()     & Fecha y hora de registro. \\
\bottomrule
\end{tabular}
\end{center}

% ----------------------------------------------------------
\subsection{Tabla: PROGRAMA}

\begin{center}
\begin{tabular}{p{4cm} p{2.8cm} p{2.8cm} p{4cm}}
\toprule
\textbf{Atributo} & \textbf{Tipo} & \textbf{Restricción} & \textbf{Descripción} \\
\midrule
programa\_id               & INT          & CP, NO NULO   & Identificador único del programa. \\
programa\_universidad\_id  & INT          & CE, NO NULO   & Universidad a la que pertenece. \\
programa\_nombre           & VARCHAR(150) & NO NULO       & Nombre del programa curricular. \\
programa\_facultad         & VARCHAR(150) & Admite nulo   & Facultad que lo administra. \\
programa\_total\_creditos  & SMALLINT     & NO NULO       & Créditos totales para graduarse. \\
programa\_created\_at      & TIMESTAMP    & DEF NOW()     & Fecha de registro. \\
\bottomrule
\end{tabular}
\end{center}

% ----------------------------------------------------------
\subsection{Tabla: TIPOLOGIA}

\begin{center}
\begin{tabular}{p{4.2cm} p{2.8cm} p{2.8cm} p{3.8cm}}
\toprule
\textbf{Atributo} & \textbf{Tipo} & \textbf{Restricción} & \textbf{Descripción} \\
\midrule
tipologia\_id                    & INT       & CP, NO NULO   & Identificador único. \\
tipologia\_programa\_id          & INT       & CE, NO NULO   & Programa al que pertenece. \\
tipologia\_nombre                & VARCHAR(100) & NO NULO    & Nombre (ej.\ Disciplinar Obligatoria). \\
tipologia\_creditos\_requeridos  & SMALLINT  & NO NULO       & Créditos mínimos requeridos en esta tipología. \\
tipologia\_cuenta\_promedio      & TINYINT   & NO NULO, DEF 1 & 1 = cuenta en el promedio; 0 = no cuenta. \\
tipologia\_created\_at           & TIMESTAMP & DEF NOW()     & Fecha de registro. \\
\bottomrule
\end{tabular}
\end{center}

% ----------------------------------------------------------
\subsection{Tabla: MATERIA}

\begin{center}
\begin{tabular}{p{4.2cm} p{2.8cm} p{2.8cm} p{3.8cm}}
\toprule
\textbf{Atributo} & \textbf{Tipo} & \textbf{Restricción} & \textbf{Descripción} \\
\midrule
materia\_id                       & INT          & CP, NO NULO          & Identificador único. \\
materia\_programa\_id             & INT          & CE, NO NULO          & Programa al que pertenece. \\
materia\_tipologia\_id            & INT          & CE, NO NULO          & Tipología de la materia. \\
materia\_codigo                   & VARCHAR(30)  & Admite nulo, Único   & Código libre por universidad (ej.\ 1000003-B). \\
materia\_nombre                   & VARCHAR(150) & NO NULO              & Nombre oficial de la asignatura. \\
materia\_creditos                 & TINYINT      & NO NULO, DEF 3       & Número de créditos académicos. \\
materia\_nota\_minima\_aprobacion & DECIMAL(3,1) & NO NULO, DEF 3.0     & Nota mínima para aprobar. \\
materia\_semestre\_sugerido       & TINYINT      & Admite nulo          & Semestre recomendado en el roadmap. \\
materia\_created\_at              & TIMESTAMP    & DEF NOW()            & Fecha de registro. \\
\bottomrule
\end{tabular}
\end{center}

% ----------------------------------------------------------
\subsection{Tabla: USUARIO}

\begin{center}
\begin{tabular}{p{4.2cm} p{2.8cm} p{2.8cm} p{3.8cm}}
\toprule
\textbf{Atributo} & \textbf{Tipo} & \textbf{Restricción} & \textbf{Descripción} \\
\midrule
usuario\_id                  & INT          & CP, NO NULO      & Identificador único del estudiante. \\
usuario\_programa\_id        & INT          & CE, NO NULO      & Programa en el que está matriculado. \\
usuario\_nombre              & VARCHAR(100) & NO NULO          & Nombre(s) del estudiante. \\
usuario\_apellido            & VARCHAR(100) & NO NULO          & Apellido(s) del estudiante. \\
usuario\_email               & VARCHAR(150) & NO NULO, Único   & Correo electrónico de acceso. \\
usuario\_password\_hash      & VARCHAR(255) & NO NULO          & Contraseña cifrada (hash). \\
usuario\_codigo\_estudiantil & VARCHAR(20)  & Admite nulo      & Código estudiantil institucional. \\
usuario\_avatar\_url         & VARCHAR(255) & Admite nulo      & URL de foto de perfil. \\
usuario\_created\_at         & TIMESTAMP    & DEF NOW()        & Fecha de registro en la plataforma. \\
\bottomrule
\end{tabular}
\end{center}

% ----------------------------------------------------------
\subsection{Tabla: SEMESTRE}

\begin{center}
\begin{tabular}{p{4.2cm} p{2.8cm} p{2.8cm} p{3.8cm}}
\toprule
\textbf{Atributo} & \textbf{Tipo} & \textbf{Restricción} & \textbf{Descripción} \\
\midrule
semestre\_id          & INT      & CP, NO NULO              & Identificador único. \\
semestre\_usuario\_id & INT      & CE, NO NULO              & Estudiante al que pertenece. \\
semestre\_numero      & TINYINT  & NO NULO                  & Número de semestre del estudiante (1, 2, 3\ldots). \\
semestre\_year        & YEAR     & NO NULO                  & Año académico (ej.\ 2025). \\
semestre\_periodo     & TINYINT  & NO NULO, CHECK (1 o 2)   & Período: 1 = primer semestre, 2 = segundo. \\
semestre\_created\_at & TIMESTAMP & DEF NOW()               & Fecha de creación del registro. \\
\bottomrule
\end{tabular}
\end{center}

% ----------------------------------------------------------
\subsection{Tabla: MATERIA\_USUARIO}

\begin{center}
\begin{tabular}{p{4.6cm} p{2.8cm} p{2.8cm} p{3.4cm}}
\toprule
\textbf{Atributo} & \textbf{Tipo} & \textbf{Restricción} & \textbf{Descripción} \\
\midrule
materia\_usuario\_id           & INT       & CP, NO NULO       & Identificador único de la inscripción. \\
materia\_usuario\_usuario\_id  & INT       & CE, NO NULO       & Estudiante que cursa la materia. \\
materia\_usuario\_materia\_id  & INT       & CE, NO NULO       & Materia inscrita. \\
materia\_usuario\_semestre\_id & INT       & CE, NO NULO       & Semestre en que se cursa. \\
materia\_usuario\_estado       & ENUM      & NO NULO, DEF en\_curso & Estado: en\_curso / aprobada / reprobada / retirada. \\
materia\_usuario\_created\_at  & TIMESTAMP & DEF NOW()         & Fecha de registro de la inscripción. \\
\bottomrule
\end{tabular}
\end{center}

% ----------------------------------------------------------
\subsection{Tabla: COMPONENTE}

\begin{center}
\begin{tabular}{p{4.2cm} p{2.8cm} p{2.8cm} p{3.8cm}}
\toprule
\textbf{Atributo} & \textbf{Tipo} & \textbf{Restricción} & \textbf{Descripción} \\
\midrule
componente\_id                   & INT          & CP, NO NULO   & Identificador único del componente. \\
componente\_materia\_usuario\_id & INT          & CE, NO NULO   & Inscripción a la que pertenece. \\
componente\_nombre               & VARCHAR(100) & NO NULO       & Nombre del componente (ej.\ Parcial 01, Talleres). \\
componente\_porcentaje           & DECIMAL(5,2) & NO NULO       & Peso en la nota final (0.00--100.00). \\
componente\_nota\_minima         & DECIMAL(3,1) & Admite nulo   & Nota mínima para superar el componente. \\
componente\_orden                & TINYINT      & NO NULO, DEF 1 & Orden de visualización. \\
componente\_created\_at          & TIMESTAMP    & DEF NOW()     & Fecha de creación. \\
\bottomrule
\end{tabular}
\end{center}

% ----------------------------------------------------------
\subsection{Tabla: NOTA}

\begin{center}
\begin{tabular}{p{4.2cm} p{2.8cm} p{2.8cm} p{3.8cm}}
\toprule
\textbf{Atributo} & \textbf{Tipo} & \textbf{Restricción} & \textbf{Descripción} \\
\midrule
nota\_id              & INT          & CP, NO NULO              & Identificador único de la nota. \\
nota\_componente\_id  & INT          & CE, NO NULO              & Componente al que pertenece la nota. \\
nota\_nombre          & VARCHAR(100) & NO NULO                  & Identificador de la entrega (ej.\ Taller 1). \\
nota\_valor           & DECIMAL(3,1) & NO NULO, CHECK 0.0--5.0  & Calificación obtenida en escala 0.0--5.0. \\
nota\_fecha\_registro & DATE         & Admite nulo              & Fecha en que se registró la nota. \\
nota\_observacion     & VARCHAR(255) & Admite nulo              & Comentario opcional del estudiante. \\
nota\_created\_at     & TIMESTAMP    & DEF NOW()                & Fecha de creación del registro. \\
\bottomrule
\end{tabular}
\end{center}

% ==========================================================
\section{Conclusiones}
\label{sec:conc}
% ==========================================================

El diseño conceptual de \textbf{UniGrades} demuestra que un modelo relacional de
9 entidades es suficiente para capturar la complejidad del progreso académico
universitario de forma agnóstica a la institución. La separación entre estructura
institucional (universidad, programa, tipología, materia) y datos personales del
estudiante (usuario, semestre, inscripción, componente, nota) garantiza escalabilidad
y reutilización de la información curricular entre múltiples usuarios.

La decisión de permitir múltiples notas por componente (relación 1:N entre
\texttt{COMPONENTE} y \texttt{NOTA}) representa la principal ventaja diferencial
del sistema frente a soluciones similares, ya que modela de forma natural los
componentes acumulativos cuya cantidad de entregas no se conoce de antemano.

Las siete reglas de negocio identificadas cubren los requisitos mínimos del modelo
y establecen las bases para las restricciones de integridad que se implementarán en
el diseño físico durante la Sesión 2.

\end{document}
